\section{团队文化与培训}
\subsection{文化落地策略}
Graham 强调：Agent 没有强力表现的时候，人们会怀疑其价值；因此必须在团队内部传播“Agent + 人类”的文化，举例说明 Agent 负责重复任务，而人类保留最终审查。每个 sprint 复盘时都要展示 Agent 成果。

\begin{enumerate}
    \item 每周分享 Agent 成功案例（如 Agent 生成 diff 并通过 3 次测试）。
    \item 制定 “Agent Etiquette” 文档，明确什么时候可以调用 Agent，哪里必须人工复盘。
    \item 设立 Mentor 制度，让新成员在熟悉 Agent 工具前由经验工程师带领。
\end{enumerate}

\subsection{培训与指南}
培训材料应包含基本术语（Plan-Act-Verify/Observability Stream）、常用 scripts 与风险提示。培训结束后可以安排实战演练：让学员在 sandbox 中让 Agent 处理一个 Issue，并在旁边观察、记录干预点。

\begin{warningbox}{培训常见挑战}
新成员容易过度依赖 Agent，建议在实战练习中设置“人工复审”门槛，确保他们熟悉如何验证 Agent 的输出。
\end{warningbox}

\begin{knowledgebox}{知识传承}
把每次人机协作的流程记录成文档与视频，便于跨团队共享经验，减少“Agent 只能某个人才能用”的风险。
\end{knowledgebox}

\subsection{本章小结}
\begin{itemize}
    \item 文化、培训与实战演练共同支撑 Agent 的长期落地。
    \item 明确的 etiquette 与复盘机制可以消除对 Agent 的不信任。
\end{itemize}

